% Generated draft: requires mathematical and source review.
\documentclass[11pt]{article}
\usepackage[margin=25mm]{geometry}
\usepackage{fontspec}
\setmainfont{Latin Modern Roman}
\usepackage{amsmath,amssymb,mathtools,unicode-math}
\setmathfont{Latin Modern Math}
\usepackage{enumitem,needspace,xurl,hyperref}
\setlength{\parindent}{0pt}
\setlength{\parskip}{5pt}
\setlength{\emergencystretch}{3em}
\newcommand{\N}{\mathbb{N}}
\newcommand{\Z}{\mathbb{Z}}
\newcommand{\Q}{\mathbb{Q}}
\newcommand{\R}{\mathbb{R}}
\newcommand{\C}{\mathbb{C}}
\newcommand{\F}{\mathbb{F}}
\newcommand{\A}{\mathbb{A}}
\newcommand{\PP}{\mathbb{P}}
\newcommand{\E}{\mathbb{E}}
\newcommand{\eps}{\varepsilon}
\newcommand{\dd}{\,\mathrm{d}}
\begin{document}
\section*{UnsolvedMath ID: 800039}
\textbf{Source problem code:} catalogue entry
\subsection{Definitions and mathematical statement}
A semigroup is a nonempty set equipped with an associative binary operation. A semigroup \(S\) is periodic if, for every \(s\in S\), there exist integers \(1\leq m<n\) such that \(s^m=s^n\).

The supplied record does not define an \emph{inclusive variety}, an \emph{identical inclusion}, or an \emph{elementary identical inclusion}. Consequently, the following is a conditional formulation of the requested classifications, pending those definitions. Let \(\mathcal I\) denote the source's class of inclusive varieties of semigroups, ordered by containment of their classes of members. For \(\mathcal V\in\mathcal I\), let \(D(\mathcal V)\) mean that \(\mathcal V\) can be defined by elementary identical inclusions, and let \(E(\mathcal V)\) mean that \(\mathcal V\) is elementary in the source's terminology. These predicates remain undefined by the supplied evidence; in particular, their equivalence is not assumed. Write \(\mathcal W\subsetneq\mathcal V\) for proper containment, and put
\[
\mathcal I_{\mathrm{per}}=\{\mathcal V\in\mathcal I: \text{every member of }\mathcal V\text{ is periodic}\}.
\]
The interpretation of the source's periodic restriction by this displayed definition requires confirmation.

G. Mashevitzky asks for the following classifications.
\begin{enumerate}
\item Find all minimal and all maximal elements of
\[
\{\mathcal V\in\mathcal I:\neg D(\mathcal V)\}.
\]
Explicitly, the minimal elements satisfy
\[
\neg D(\mathcal V)\quad\text{and}\quad
\forall\mathcal W\in\mathcal I\,
\bigl(\mathcal W\subsetneq\mathcal V\Longrightarrow D(\mathcal W)\bigr),
\]
and the maximal elements satisfy
\[
\neg D(\mathcal V)\quad\text{and}\quad
\forall\mathcal W\in\mathcal I\,
\bigl(\mathcal V\subsetneq\mathcal W\Longrightarrow D(\mathcal W)\bigr).
\]
\item Describe all \(\mathcal V\in\mathcal I\) such that
\[
E(\mathcal V)\quad\text{and}\quad
\forall\mathcal W\in\mathcal I\,
\bigl(\mathcal W\subseteq\mathcal V\Longrightarrow E(\mathcal W)\bigr).
\]
\item Describe all \(\mathcal V\in\mathcal I_{\mathrm{per}}\) such that
\[
\neg E(\mathcal V)\quad\text{and}\quad
\forall\mathcal W\in\mathcal I_{\mathrm{per}}\,
\bigl(\mathcal V\subseteq\mathcal W\Longrightarrow\neg E(\mathcal W)\bigr).
\]
Here the quantification over supervarieties is restricted to inclusive varieties consisting of periodic semigroups.
\end{enumerate}
The supplied statement additionally reports that all minimal inclusive varieties of Abelian groups that cannot be defined by elementary identical inclusions are described in reference [M3]. This is a supplied background assertion, not an independently verified result; the bibliographic identity of [M3] is missing.
\subsection{Short English statement}
The problem asks for classifications of minimal and maximal inclusive varieties not definable by elementary identical inclusions, elementary inclusive varieties whose inclusive subvarieties are all elementary, and nonelementary inclusive varieties of periodic semigroups whose periodic inclusive supervarieties are all nonelementary. The technical meanings of inclusive variety and elementary are missing, so the displayed formulation remains conditional.
\subsection{Sources}
\begin{itemize}
\item Local UnsolvedMath registry, permanent ID \(800039\), source problem code \texttt{catalogue entry}, title \emph{The Lyapin's Notebook --- Problem 41}. The supplied registry labels the problem open; that label is not an independently verified literature assessment.
\item Bershadsky, Kublanovsky, and Mashevitzky, \emph{The Lyapin's Notebook: Unsolved Problems in Semigroup Theory} (2026), Problem 41, attributed in the statement to G. Mashevitzky. Supplied source URL: \url{https://arxiv.org/abs/2604.04763}. The registry's proposer field lists Bershadsky, Kublanovsky, and Mashevitzky and gives proposed year 2026.
\item Supplied provenance records access on 2026-10-16 and extraction by \texttt{arxiv\_latex}. These are supplied provenance notes; no independent source visit or verification is asserted here.
\item Reference [M3], cited in the supplied statement for the classification concerning Abelian groups. No bibliographic details or separate URL for [M3] were supplied.
\end{itemize}
\end{document}
